\section{Carrier access}\label{supp:access}
These calculations support the distinction between applied field and local
polarization in \S\ref{sec:khz} and its kHz demonstrations.
\subsection{Compartment access and the passive-cable estimate}\label{app:access}
The passive compartment equation $\mathsf C_m\dot{\mathbf V}+\mathsf G_m\mathbf V=\mathbf I^0-\mathbf p_EE$
\cite{ruffini2021nmm2} gives the voltage perturbation
$\delta\mathbf V(\omega)=-(\mathsf G_m+i\omega\mathsf C_m)^{-1}\mathbf p_EE(\omega)$. In SI units, $\mathbf V$ contains the compartment membrane potentials (V), $\mathsf C_m$ and
$\mathsf G_m$ are the capacitance (F) and conductance (S) matrices, $\mathbf I^0$ is the
endogenous current (A), and $\mathbf p_E$ is the field-to-current coupling
(A per V/m).

Geometry can also alter the relevant response mode. For a passive cable, the
magnitude of the frequency-dependent space constant is
\begin{equation}
\ell(\omega)=\ell_0(1+\omega^2\tau_m^2)^{-1/4},
\label{eq:lambda}
\end{equation}
where $\ell_0$ is its static value. At $1$~kHz and $\tau_m=16$~ms, this length
is about one tenth of its static value, so compartments that are strongly coupled
at low frequency can respond differently to the carrier. A somatic charging time
alone therefore does not determine every extracellular-field response mode.
For a sufficiently long passive cable with a sealed end in a uniform axial field,
the end polarization is proportional to this space constant. Each carrier then
acquires the factor $(1+\omega^2\tau_m^2)^{-1/4}$, so the quadratic beat is
attenuated by $(1+\omega_c^2\tau_m^2)^{-1/2}$, about $10^{-2}$ at $1$~kHz and
$10^{-3}$ at $10$~kHz for $16$~ms. These differ from the point-membrane estimates;
soma loading and finite cable length further change the transfer.
Thus cable geometry can improve carrier access without changing the somatic
membrane time.

\subsection{Measured responses and candidate access sites}
Cohen et al.'s periaxonal circuit provides a reason to distinguish axonal from
somatic access \cite{cohen2020periaxonal}. A conducting submyelin space, about
$12$~nm wide and partly sealed at the paranodes, separates the voltage across
myelin from that across the axonal membrane. Their minimum fitted axonal decay
was approximately $1$~ms. This is a response to injected current in a distributed
cable, rather than a measured transfer from a TI field. The AIS is another candidate
site, with a high sodium-channel density \cite{kole2008ais}; this property supports
nonlinear sensitivity but does not fix its charging time.

\ifjne
\subsection{Access assumptions in related models}
Short access times motivate the model but remain preparation-dependent.
The $1$ and $0.4$~ms cases of \cite{esmaeilpour2021} precede our effective-time
hypothesis. Luff et al.'s pulse-width-modulated TI model uses $1$~ms within a
$1$--$100$~ms passive-time sweep \cite{luff2024}. The $0.7~\mu$s charging
measured by Semenov et al.\ under strong microsecond field pulses
\cite{semenov2025} concerns a different response mode from intracellular
current charging or periaxonal conduction. A $10$--$100$-fold reduction in
QIF membrane time is therefore a sensitivity hypothesis, whose population
consequence follows from the slow-synapse limit.

\fi
\setcounter{figure}{4}
\begin{figure}[t!]
\centering
\includegraphics[width=0.74\textwidth]{fig_khz.pdf}
\caption{\textbf{Carrier attenuation depends on the assumed access time (Jansen--Rit, open loop).} Open-loop
demodulated response (normalized to the $\tau\!\to\!0$ value) vs.\ carrier
frequency up to $10$~kHz. With direct coupling ($\tau\!\to\!0$) the response is
flat---carrier-independent. A slow somatic membrane ($\tau_m\!\approx\!16$~ms)
rolls off as $1/f_c^2$ (dashed guide); the assumed fast response
($\tau=0.2$~ms) retains about $0.14$ of the direct response at $2$~kHz and $0.04$
at $3.9$~kHz. These are sensitivity cases, not measurements of anatomical
compartments.}
\label{fig:khz}
\end{figure}

\setcounter{figure}{5}
\begin{figure}[t!]
\centering
\includegraphics[width=0.95\textwidth]{fig_khz_twotone.pdf}
\caption{\textbf{Two-tone TI at a kHz carrier (Jansen--Rit, assumed fast access).}
Two equal tones at $f_1=1994.543$ and $f_2=2005.457$~Hz, beating at
$10.914$~Hz, pass through the assumed low-pass ($\tau=0.2$~ms) into the recurrent column.
\emph{(a)} The input amplitude spectrum has lines only at $f_1$ and $f_2$; in the
$8$--$14$~Hz band its maximum is $2.2\times10^{-13}$~mV.
\emph{(b)} The output carries a demodulated line at the beat, $10.9$~Hz.}
\label{fig:khz_twotone}
\end{figure}

\setcounter{figure}{8}
\begin{figure}[t!]
\centering
\includegraphics[width=\textwidth]{fig_khz_direct.pdf}
\caption{\textbf{Demodulation at a kHz carrier (Jansen--Rit, closed loop).}
The applied field passes through an explicit membrane filter before the sigmoid.
\emph{(a)} Closed-loop demodulated $A_\Omega$ at $\Delta f=f_0$ vs.\ carrier $f_c$ for
three access models: direct ($\tau\!\to\!0$, carrier-independent), fast
($\tau=0.2$~ms), and somatic ($\tau_m=16$~ms; $1/f_c^2$ roll-off).
The fast time is assumed rather than assigned to a measured compartment.
\emph{(b)} Output spectrum for a
$2$~kHz carrier delivered through the fast element: an alpha line at $f_0$,
synthesized by the column although the input carries no alpha power.}
\label{fig:khz_direct}
\end{figure}

